% !TEX program = pdflatex
% !TEX encoding = UTF-8 Unicode
\documentclass[12pt,a4paper]{article}
\usepackage[utf8]{vietnam}
\usepackage{amsmath,amssymb}
\usepackage{tabularx}
\usepackage{tikz}
\usepackage{xcolor}
\usepackage[top=0.45cm, bottom=0.45cm, left=0.6cm, right=0.6cm]{geometry}
\usepackage[dethi]{ex_test}

\setlength{\parskip}{0pt}
\setlength{\parindent}{0pt}

% Dòng kẻ chấm mờ đều đặn chuẩn viết tay (khoảng cách 19.5pt)
\newcommand{\dongke}[1]{%
  \begingroup%
  \setlength{\parskip}{0pt}%
  \baselineskip=19.5pt%
  \vspace{1mm}%
  \foreach \i in {1,...,#1}{%
    \noindent{\color{gray!60}\dotfill}\par%
  }%
  \vspace{0.8mm}%
  \endgroup%
}

% Định nghĩa ô bài tập tự luận trong bảng
\newcommand{\oBaiTap}[3][5]{%
  \vspace{0.8mm}%
  \noindent\textbf{#2} #3\par%
  \vspace{0.8mm}%
  \noindent\textit{Bài giải:}%
  \dongke{#1}%
}

% Câu trắc nghiệm 4 đáp án dàn hàng ngang
\newcommand{\cauTN}[6]{%
  \noindent\textbf{Câu #1.} #2\par\vspace{1mm}%
  \noindent%
  \makebox[0.25\linewidth][l]{\textbf{A.} #3}%
  \makebox[0.25\linewidth][l]{\textbf{B.} #4}%
  \makebox[0.25\linewidth][l]{\textbf{C.} #5}%
  \makebox[0.25\linewidth][l]{\textbf{D.} #6}\par\vspace{1.5mm}%
}

% Câu trắc nghiệm 2 dòng (mỗi dòng 2 đáp án, chống lỗi đè chữ khi biểu thức dài)
\newcommand{\cauTNHaiDong}[6]{%
  \noindent\textbf{Câu #1.} #2\par\vspace{1mm}%
  \noindent%
  \makebox[0.5\linewidth][l]{\textbf{A.} #3}%
  \makebox[0.5\linewidth][l]{\textbf{B.} #4}\par\vspace{1mm}%
  \noindent%
  \makebox[0.5\linewidth][l]{\textbf{C.} #5}%
  \makebox[0.5\linewidth][l]{\textbf{D.} #6}\par\vspace{1.5mm}%
}

\begin{document}
\pagestyle{empty}

% #############################################################################
% ========================= TOÁN 8 - TUẦN 04 - BUỔI 02: MÃ ĐỀ 101 ============
% ========================= MẶT TRƯỚC (TRANG 1) ===============================
% #############################################################################

\noindent
\begin{tabularx}{\linewidth}{|X|c|}
\hline
\textbf{PHIẾU KIỂM TRA ĐỊNH KỲ TOÁN 8 - TUẦN 04 - BUỔI 02} & \textbf{Mã đề: 101} \\
\textit{Chủ đề: Lập phương của một tổng hoặc một hiệu} & \textit{Thời gian: 45 phút} \\
\hline
\end{tabularx}

\vspace{1mm}

\noindent
\begin{tabularx}{\linewidth}{|X|}
\hline
\rule[-2.5mm]{0pt}{8mm}\textbf{Họ và tên học sinh:} \dotfill \quad \textbf{Mã HS:} \makebox[3.5cm]{\dotfill} \\
\hline
\end{tabularx}

\vspace{1mm}

\noindent
\begin{tabularx}{\linewidth}{|p{3.2cm}|X|}
\hline
\centering\textbf{ĐIỂM SỐ} & \centering\textbf{LỜI NHẬN XÉT CỦA THẦY CÔ} \tabularnewline
\hline
\rule{0pt}{1.3cm} & \tabularnewline
\hline
\end{tabularx}

\vspace{1.2mm}

\noindent\textbf{BẢNG ĐIỀN ĐÁP ÁN TRẮC NGHIỆM (Mỗi câu đúng 0,5 điểm):}
\vspace{0.8mm}

\noindent
\renewcommand{\arraystretch}{1.35}
\begin{tabularx}{\linewidth}{|c|*{5}{>{\centering\arraybackslash}X|}}
\hline
\textbf{Câu} & \textbf{1} & \textbf{2} & \textbf{3} & \textbf{4} & \textbf{5} \\
\hline
\textbf{Đáp án} & & & & & \\
\hline
\end{tabularx}
\renewcommand{\arraystretch}{1.0}

\vspace{1.5mm}

\noindent\textbf{\large PHẦN I. TRẮC NGHIỆM KHÁCH QUAN (2,5 điểm)}
\vspace{1mm}

\cauTNHaiDong{1}{Hằng đẳng thức lập phương của một tổng $(A+B)^3$ là:}{$A^3 + 3A^2B + 3AB^2 + B^3$}{$A^3 - 3A^2B + 3AB^2 - B^3$}{$A^3 + B^3$}{$A^3 + 3A^2B + B^3$}

\cauTN{2}{Khai triển hằng đẳng thức $(x + 2)^3$ ta được kết quả là:}{$x^3 + 8$}{$x^3 + 3x^2 + 6x + 8$}{$x^3 + 6x^2 + 12x + 8$}{$x^3 + 6x^2 + 6x + 8$}

\cauTN{3}{Thu gọn đa thức $x^3 - 3x^2 + 3x - 1$ về dạng lập phương của một hiệu:}{$(x - 3)^3$}{$(x + 1)^3$}{$(x - 1)^3$}{$(1 - x)^2$}

\cauTN{4}{Hệ số của hạng tử $x^2$ trong khai triển biểu thức $(2x - 1)^3$ là:}{$-12$}{$12$}{$-6$}{$8$}

\cauTN{5}{Khẳng định nào dưới đây là \textbf{đúng}?}{$(A - B)^3 = -(B - A)^3$}{$(A - B)^3 = (B - A)^3$}{$(A + B)^3 = A^3 + B^3$}{$(A - B)^3 = A^3 - B^3$}

\vspace{1mm}

\noindent\textbf{\large PHẦN II. TỰ LUẬN (7,5 điểm)}

\noindent\textbf{Câu 1 (2,5 điểm). Rút gọn các biểu thức sau:}
\vspace{0.8mm}

\noindent
\begin{tabularx}{\linewidth}{|X|X|}
\hline
\oBaiTap[6]{a)}{$A = (x + 2)^3 + (x - 2)^3$} & 
\oBaiTap[6]{b)}{$B = (x + 1)^3 - (x - 1)^3 - 6x^2$} \\
\hline
\oBaiTap[6]{c)}{$C = (x - 1)^3 - x(x + 2)(x - 2) + 3x^2$} & 
\oBaiTap[6]{d)}{$D = (x + 1)^3 - x^2(x + 3) - 3x$} \\
\hline
\end{tabularx}

\newpage
% #############################################################################
% ========================= TOÁN 8 - TUẦN 04 - BUỔI 02: MÃ ĐỀ 101 ============
% ========================= MẶT SAU (TRANG 2) =================================
% #############################################################################

\noindent
\begin{tabularx}{\linewidth}{|X|c|}
\hline
\textbf{PHẦN II. TỰ LUẬN (Tiếp theo) - TOÁN 8 TUẦN 04 BUỔI 02} & \textbf{Mã đề: 101} \\
\hline
\end{tabularx}

\vspace{1mm}

\noindent
\begin{tabularx}{\linewidth}{|X|}
\hline
\rule[-2mm]{0pt}{7mm}\textbf{Họ và tên học sinh:} \dotfill \quad \textbf{Mã HS:} \makebox[3.5cm]{\dotfill} \\
\hline
\end{tabularx}

\vspace{1.5mm}

\noindent\textbf{\large Câu 2 (2,5 điểm). Rút gọn và tìm $x$:}
\vspace{0.8mm}

\noindent
\begin{tabularx}{\linewidth}{|X|X|}
\hline
\oBaiTap[10]{a)}{Rút gọn và tính giá trị biểu thức:\par
$P = x^3 - 3x^2 + 3x - 1$\par
tại $x = 101$.} & 
\oBaiTap[10]{b)}{Tìm số thực $x$, biết:\par
$(x + 1)^3 - x^2(x + 3) = 19$.} \\
\hline
\end{tabularx}

\vspace{1.8mm}

\noindent\textbf{\large Câu 3 (2,5 điểm). Bài toán thực tế và nâng cao:}
\vspace{0.8mm}

\noindent
\begin{tabularx}{\linewidth}{|X|}
\hline
\vspace{0.5mm}
\textbf{a) (1,5 điểm)} Một bể chứa nước sinh hoạt bằng inox có dạng hình lập phương cạnh $x$ (mét). Do nhu cầu mở rộng, gia đình làm một bể mới có cạnh tăng thêm $0{,}5\text{ m}$.\par
1) Hãy viết biểu thức biểu thị thể tích bể nước mới theo $x$. Viết biểu thức biểu thị lượng thể tích nước tăng thêm $\Delta V$ của bể mới so với bể cũ theo $x$.\par
2) Cho biết cạnh bể ban đầu là $x = 1{,}5\text{ m}$. Hãy tính chính xác lượng thể tích nước tăng thêm $\Delta V$ (theo $\text{m}^3$ và đổi ra lít nước, biết $1\text{ m}^3 = 1000$ lít).\par
\vspace{0.4mm}
\noindent\textit{Bài giải:}\dongke{7}
\vspace{1mm}
\textbf{b) (1,0 điểm)} Chứng minh rằng giá trị của biểu thức sau không phụ thuộc vào biến $x$:\par
$M = (x + 2)^3 - (x - 2)^3 - 12x^2$.\par
\vspace{0.4mm}
\noindent\textit{Bài giải:}\dongke{6}
\\ \hline
\end{tabularx}

\newpage
% #############################################################################
% ========================= TOÁN 8 - TUẦN 04 - BUỔI 02: MÃ ĐỀ 102 ============
% ========================= MẶT TRƯỚC (TRANG 3) ===============================
% #############################################################################

\noindent
\begin{tabularx}{\linewidth}{|X|c|}
\hline
\textbf{PHIẾU KIỂM TRA ĐỊNH KỲ TOÁN 8 - TUẦN 04 - BUỔI 02} & \textbf{Mã đề: 102} \\
\textit{Chủ đề: Lập phương của một tổng hoặc một hiệu} & \textit{Thời gian: 45 phút} \\
\hline
\end{tabularx}

\vspace{1mm}

\noindent
\begin{tabularx}{\linewidth}{|X|}
\hline
\rule[-2.5mm]{0pt}{8mm}\textbf{Họ và tên học sinh:} \dotfill \quad \textbf{Mã HS:} \makebox[3.5cm]{\dotfill} \\
\hline
\end{tabularx}

\vspace{1mm}

\noindent
\begin{tabularx}{\linewidth}{|p{3.2cm}|X|}
\hline
\centering\textbf{ĐIỂM SỐ} & \centering\textbf{LỜI NHẬN XÉT CỦA THẦY CÔ} \tabularnewline
\hline
\rule{0pt}{1.3cm} & \tabularnewline
\hline
\end{tabularx}

\vspace{1.2mm}

\noindent\textbf{BẢNG ĐIỀN ĐÁP ÁN TRẮC NGHIỆM (Mỗi câu đúng 0,5 điểm):}
\vspace{0.8mm}

\noindent
\renewcommand{\arraystretch}{1.35}
\begin{tabularx}{\linewidth}{|c|*{5}{>{\centering\arraybackslash}X|}}
\hline
\textbf{Câu} & \textbf{1} & \textbf{2} & \textbf{3} & \textbf{4} & \textbf{5} \\
\hline
\textbf{Đáp án} & & & & & \\
\hline
\end{tabularx}
\renewcommand{\arraystretch}{1.0}

\vspace{1.5mm}

\noindent\textbf{\large PHẦN I. TRẮC NGHIỆM KHÁCH QUAN (2,5 điểm)}
\vspace{1mm}

\cauTNHaiDong{1}{Hằng đẳng thức lập phương của một hiệu $(A-B)^3$ là:}{$A^3 - 3A^2B + 3AB^2 - B^3$}{$A^3 + 3A^2B + 3AB^2 + B^3$}{$A^3 - B^3$}{$A^3 - 3A^2B - B^3$}

\cauTN{2}{Khai triển hằng đẳng thức $(x - 2)^3$ ta được kết quả là:}{$x^3 - 8$}{$x^3 - 6x^2 + 12x - 8$}{$x^3 - 6x^2 + 6x - 8$}{$x^3 - 3x^2 + 6x - 8$}

\cauTN{3}{Thu gọn đa thức $x^3 + 6x^2 + 12x + 8$ về dạng lập phương của một tổng:}{$(x + 8)^3$}{$(x + 4)^3$}{$(x + 2)^3$}{$(x + 6)^3$}

\cauTN{4}{Hệ số của hạng tử $x$ trong khai triển biểu thức $(x + 3)^3$ là:}{$27$}{$9$}{$3$}{$1$}

\cauTN{5}{Khẳng định nào dưới đây là \textbf{sai}?}{$(x - y)^3 = -(y - x)^3$}{$(x + y)^3 = (y + x)^3$}{$(x - y)^3 = x^3 - 3x^2y + 3xy^2 - y^3$}{$(x - y)^3 = (y - x)^3$}

\vspace{1mm}

\noindent\textbf{\large PHẦN II. TỰ LUẬN (7,5 điểm)}

\noindent\textbf{Câu 1 (2,5 điểm). Rút gọn các biểu thức sau:}
\vspace{0.8mm}

\noindent
\begin{tabularx}{\linewidth}{|X|X|}
\hline
\oBaiTap[6]{a)}{$A = (x + 3)^3 + (x - 3)^3$} & 
\oBaiTap[6]{b)}{$B = (x + 2)^3 - (x - 2)^3 - 12x^2$} \\
\hline
\oBaiTap[6]{c)}{$C = (x - 2)^3 - x(x + 3)(x - 3) + 6x^2$} & 
\oBaiTap[6]{d)}{$D = (x + 2)^3 - x^2(x + 6) - 12x$} \\
\hline
\end{tabularx}

\newpage
% #############################################################################
% ========================= TOÁN 8 - TUẦN 04 - BUỔI 02: MÃ ĐỀ 102 ============
% ========================= MẶT SAU (TRANG 4) =================================
% #############################################################################

\noindent
\begin{tabularx}{\linewidth}{|X|c|}
\hline
\textbf{PHẦN II. TỰ LUẬN (Tiếp theo) - TOÁN 8 TUẦN 04 BUỔI 02} & \textbf{Mã đề: 102} \\
\hline
\end{tabularx}

\vspace{1mm}

\noindent
\begin{tabularx}{\linewidth}{|X|}
\hline
\rule[-2mm]{0pt}{7mm}\textbf{Họ và tên học sinh:} \dotfill \quad \textbf{Mã HS:} \makebox[3.5cm]{\dotfill} \\
\hline
\end{tabularx}

\vspace{1.5mm}

\noindent\textbf{\large Câu 2 (2,5 điểm). Rút gọn và tìm $x$:}
\vspace{0.8mm}

\noindent
\begin{tabularx}{\linewidth}{|X|X|}
\hline
\oBaiTap[10]{a)}{Rút gọn và tính giá trị biểu thức:\par
$P = x^3 - 6x^2 + 12x - 8$\par
tại $x = 102$.} & 
\oBaiTap[10]{b)}{Tìm số thực $x$, biết:\par
$(x + 2)^3 - x^2(x + 6) = 32$.} \\
\hline
\end{tabularx}

\vspace{1.8mm}

\noindent\textbf{\large Câu 3 (2,5 điểm). Bài toán thực tế và nâng cao:}
\vspace{0.8mm}

\noindent
\begin{tabularx}{\linewidth}{|X|}
\hline
\vspace{0.5mm}
\textbf{a) (1,5 điểm)} Một bể chứa nước sinh hoạt hình lập phương cạnh $x$ (mét). Người ta nâng cấp dung tích bằng cách tăng độ dài mỗi cạnh thêm $1\text{ m}$.\par
1) Hãy viết biểu thức biểu thị thể tích bể nước mới theo $x$. Viết biểu thức biểu thị lượng thể tích nước tăng thêm $\Delta V$ của bể mới so với bể cũ theo $x$.\par
2) Cho biết cạnh bể ban đầu là $x = 1\text{ m}$. Hãy tính chính xác lượng thể tích nước tăng thêm $\Delta V$ (theo $\text{m}^3$ và đổi ra lít nước, biết $1\text{ m}^3 = 1000$ lít).\par
\vspace{0.4mm}
\noindent\textit{Bài giải:}\dongke{7}
\vspace{1mm}
\textbf{b) (1,0 điểm)} Chứng minh rằng giá trị của biểu thức sau không phụ thuộc vào biến $x$:\par
$M = (x + 3)^3 - (x - 3)^3 - 18x^2$.\par
\vspace{0.4mm}
\noindent\textit{Bài giải:}\dongke{6}
\\ \hline
\end{tabularx}

\newpage
% #############################################################################
% ========================= TOÁN 8 - TUẦN 04 - BUỔI 02: MÃ ĐỀ 103 ============
% ========================= MẶT TRƯỚC (TRANG 5) ===============================
% #############################################################################

\noindent
\begin{tabularx}{\linewidth}{|X|c|}
\hline
\textbf{PHIẾU KIỂM TRA ĐỊNH KỲ TOÁN 8 - TUẦN 04 - BUỔI 02} & \textbf{Mã đề: 103} \\
\textit{Chủ đề: Lập phương của một tổng hoặc một hiệu} & \textit{Thời gian: 45 phút} \\
\hline
\end{tabularx}

\vspace{1mm}

\noindent
\begin{tabularx}{\linewidth}{|X|}
\hline
\rule[-2.5mm]{0pt}{8mm}\textbf{Họ và tên học sinh:} \dotfill \quad \textbf{Mã HS:} \makebox[3.5cm]{\dotfill} \\
\hline
\end{tabularx}

\vspace{1mm}

\noindent
\begin{tabularx}{\linewidth}{|p{3.2cm}|X|}
\hline
\centering\textbf{ĐIỂM SỐ} & \centering\textbf{LỜI NHẬN XÉT CỦA THẦY CÔ} \tabularnewline
\hline
\rule{0pt}{1.3cm} & \tabularnewline
\hline
\end{tabularx}

\vspace{1.2mm}

\noindent\textbf{BẢNG ĐIỀN ĐÁP ÁN TRẮC NGHIỆM (Mỗi câu đúng 0,5 điểm):}
\vspace{0.8mm}

\noindent
\renewcommand{\arraystretch}{1.35}
\begin{tabularx}{\linewidth}{|c|*{5}{>{\centering\arraybackslash}X|}}
\hline
\textbf{Câu} & \textbf{1} & \textbf{2} & \textbf{3} & \textbf{4} & \textbf{5} \\
\hline
\textbf{Đáp án} & & & & & \\
\hline
\end{tabularx}
\renewcommand{\arraystretch}{1.0}

\vspace{1.5mm}

\noindent\textbf{\large PHẦN I. TRẮC NGHIỆM KHÁCH QUAN (2,5 điểm)}
\vspace{1mm}

\cauTNHaiDong{1}{Hằng đẳng thức lập phương của một tổng $(A+B)^3$ bằng:}{$A^3 + 3A^2B + 3AB^2 + B^3$}{$A^3 + B^3$}{$A^3 + 3AB + B^3$}{$A^3 + A^2B + AB^2 + B^3$}

\cauTN{2}{Khai triển hằng đẳng thức $(x + 1)^3$ ta được kết quả là:}{$x^3 + 1$}{$x^3 + 3x^2 + 1$}{$x^3 + 3x^2 + 3x + 1$}{$x^3 + x^2 + x + 1$}

\cauTN{3}{Thu gọn đa thức $x^3 - 6x^2 + 12x - 8$ về dạng lập phương của một hiệu:}{$(x - 8)^3$}{$(x - 4)^3$}{$(x - 2)^3$}{$(x - 6)^3$}

\cauTN{4}{Hệ số của hạng tử $x^2$ trong khai triển biểu thức $(3x - 1)^3$ là:}{$-27$}{$27$}{$-9$}{$9$}

\cauTN{5}{Khẳng định nào dưới đây là \textbf{sai}?}{$(a + b)^3 = a^3 + 3a^2b + 3ab^2 + b^3$}{$(a - b)^3 = a^3 - b^3$}{$-(a - b)^3 = (b - a)^3$}{$(a + b)^3 = (b + a)^3$}

\vspace{1mm}

\noindent\textbf{\large PHẦN II. TỰ LUẬN (7,5 điểm)}

\noindent\textbf{Câu 1 (2,5 điểm). Rút gọn các biểu thức sau:}
\vspace{0.8mm}

\noindent
\begin{tabularx}{\linewidth}{|X|X|}
\hline
\oBaiTap[6]{a)}{$A = (x + 1)^3 + (x - 1)^3$} & 
\oBaiTap[6]{b)}{$B = (x + 3)^3 - (x - 3)^3 - 18x^2$} \\
\hline
\oBaiTap[6]{c)}{$C = (x - 3)^3 - x(x + 2)(x - 2) + 9x^2$} & 
\oBaiTap[6]{d)}{$D = (x + 3)^3 - x^2(x + 9) - 27x$} \\
\hline
\end{tabularx}

\newpage
% #############################################################################
% ========================= TOÁN 8 - TUẦN 04 - BUỔI 02: MÃ ĐỀ 103 ============
% ========================= MẶT SAU (TRANG 6) =================================
% #############################################################################

\noindent
\begin{tabularx}{\linewidth}{|X|c|}
\hline
\textbf{PHẦN II. TỰ LUẬN (Tiếp theo) - TOÁN 8 TUẦN 04 BUỔI 02} & \textbf{Mã đề: 103} \\
\hline
\end{tabularx}

\vspace{1mm}

\noindent
\begin{tabularx}{\linewidth}{|X|}
\hline
\rule[-2mm]{0pt}{7mm}\textbf{Họ và tên học sinh:} \dotfill \quad \textbf{Mã HS:} \makebox[3.5cm]{\dotfill} \\
\hline
\end{tabularx}

\vspace{1.5mm}

\noindent\textbf{\large Câu 2 (2,5 điểm). Rút gọn và tìm $x$:}
\vspace{0.8mm}

\noindent
\begin{tabularx}{\linewidth}{|X|X|}
\hline
\oBaiTap[10]{a)}{Rút gọn và tính giá trị biểu thức:\par
$P = x^3 + 3x^2 + 3x + 1$\par
tại $x = 99$.} & 
\oBaiTap[10]{b)}{Tìm số thực $x$, biết:\par
$(x + 1)^3 - x^2(x + 3) = 16$.} \\
\hline
\end{tabularx}

\vspace{1.8mm}

\noindent\textbf{\large Câu 3 (2,5 điểm). Bài toán thực tế và nâng cao:}
\vspace{0.8mm}

\noindent
\begin{tabularx}{\linewidth}{|X|}
\hline
\vspace{0.5mm}
\textbf{a) (1,5 điểm)} Một bể nước hình lập phương có độ dài cạnh ban đầu là $x$ (mét). Người ta làm bể mới tăng cạnh thêm $0{,}5\text{ m}$.\par
1) Hãy viết biểu thức biểu thị thể tích bể nước mới theo $x$. Viết biểu thức biểu thị lượng thể tích nước tăng thêm $\Delta V$ của bể mới so với bể cũ theo $x$.\par
2) Cho biết cạnh bể ban đầu là $x = 2\text{ m}$. Hãy tính chính xác lượng thể tích nước tăng thêm $\Delta V$ (theo $\text{m}^3$ và đổi ra lít nước, biết $1\text{ m}^3 = 1000$ lít).\par
\vspace{0.4mm}
\noindent\textit{Bài giải:}\dongke{7}
\vspace{1mm}
\textbf{b) (1,0 điểm)} Chứng minh rằng giá trị của biểu thức sau không phụ thuộc vào biến $x$:\par
$M = (x + 1)^3 - (x - 1)^3 - 6x^2$.\par
\vspace{0.4mm}
\noindent\textit{Bài giải:}\dongke{6}
\\ \hline
\end{tabularx}

\newpage
% #############################################################################
% ========================= TOÁN 8 - TUẦN 04 - BUỔI 02: MÃ ĐỀ 104 ============
% ========================= MẶT TRƯỚC (TRANG 7) ===============================
% #############################################################################

\noindent
\begin{tabularx}{\linewidth}{|X|c|}
\hline
\textbf{PHIẾU KIỂM TRA ĐỊNH KỲ TOÁN 8 - TUẦN 04 - BUỔI 02} & \textbf{Mã đề: 104} \\
\textit{Chủ đề: Lập phương của một tổng hoặc một hiệu} & \textit{Thời gian: 45 phút} \\
\hline
\end{tabularx}

\vspace{1mm}

\noindent
\begin{tabularx}{\linewidth}{|X|}
\hline
\rule[-2.5mm]{0pt}{8mm}\textbf{Họ và tên học sinh:} \dotfill \quad \textbf{Mã HS:} \makebox[3.5cm]{\dotfill} \\
\hline
\end{tabularx}

\vspace{1mm}

\noindent
\begin{tabularx}{\linewidth}{|p{3.2cm}|X|}
\hline
\centering\textbf{ĐIỂM SỐ} & \centering\textbf{LỜI NHẬN XÉT CỦA THẦY CÔ} \tabularnewline
\hline
\rule{0pt}{1.3cm} & \tabularnewline
\hline
\end{tabularx}

\vspace{1.2mm}

\noindent\textbf{BẢNG ĐIỀN ĐÁP ÁN TRẮC NGHIỆM (Mỗi câu đúng 0,5 điểm):}
\vspace{0.8mm}

\noindent
\renewcommand{\arraystretch}{1.35}
\begin{tabularx}{\linewidth}{|c|*{5}{>{\centering\arraybackslash}X|}}
\hline
\textbf{Câu} & \textbf{1} & \textbf{2} & \textbf{3} & \textbf{4} & \textbf{5} \\
\hline
\textbf{Đáp án} & & & & & \\
\hline
\end{tabularx}
\renewcommand{\arraystretch}{1.0}

\vspace{1.5mm}

\noindent\textbf{\large PHẦN I. TRẮC NGHIỆM KHÁCH QUAN (2,5 điểm)}
\vspace{1mm}

\cauTNHaiDong{1}{Hằng đẳng thức lập phương của một hiệu $(A-B)^3$ bằng:}{$A^3 - 3A^2B + 3AB^2 - B^3$}{$A^3 - B^3$}{$A^3 + 3A^2B + 3AB^2 + B^3$}{$A^3 - 3AB + B^3$}

\cauTN{2}{Khai triển hằng đẳng thức $(x - 3)^3$ ta được kết quả là:}{$x^3 - 27$}{$x^3 - 9x^2 + 27x - 27$}{$x^3 - 9x^2 + 9x - 27$}{$x^3 - 3x^2 + 9x - 27$}

\cauTN{3}{Thu gọn đa thức $x^3 + 9x^2 + 27x + 27$ về dạng lập phương của một tổng:}{$(x + 9)^3$}{$(x + 3)^3$}{$(x + 27)^3$}{$(x + 1)^3$}

\cauTN{4}{Hệ số của hạng tử $x$ trong khai triển biểu thức $(2x - 1)^3$ là:}{$6$}{$-6$}{$12$}{$-12$}

\cauTN{5}{Khẳng định nào dưới đây là \textbf{sai}?}{$(a - b)^3 = -(b - a)^3$}{$(a + b)^3 = a^3 + b^3$}{$(a - b)^3 = a^3 - 3a^2b + 3ab^2 - b^3$}{$(-a - b)^3 = -(a + b)^3$}

\vspace{1mm}

\noindent\textbf{\large PHẦN II. TỰ LUẬN (7,5 điểm)}

\noindent\textbf{Câu 1 (2,5 điểm). Rút gọn các biểu thức sau:}
\vspace{0.8mm}

\noindent
\begin{tabularx}{\linewidth}{|X|X|}
\hline
\oBaiTap[6]{a)}{$A = (2x + 1)^3 + (2x - 1)^3$} & 
\oBaiTap[6]{b)}{$B = (x + 4)^3 - (x - 4)^3 - 24x^2$} \\
\hline
\oBaiTap[6]{c)}{$C = (x - 1)^3 - x(x + 3)(x - 3) + 3x^2$} & 
\oBaiTap[6]{d)}{$D = (2x + 1)^3 - 4x(2x^2 + 3x) - 1$} \\
\hline
\end{tabularx}

\newpage
% #############################################################################
% ========================= TOÁN 8 - TUẦN 04 - BUỔI 02: MÃ ĐỀ 104 ============
% ========================= MẶT SAU (TRANG 8) =================================
% #############################################################################

\noindent
\begin{tabularx}{\linewidth}{|X|c|}
\hline
\textbf{PHẦN II. TỰ LUẬN (Tiếp theo) - TOÁN 8 TUẦN 04 BUỔI 02} & \textbf{Mã đề: 104} \\
\hline
\end{tabularx}

\vspace{1mm}

\noindent
\begin{tabularx}{\linewidth}{|X|}
\hline
\rule[-2mm]{0pt}{7mm}\textbf{Họ và tên học sinh:} \dotfill \quad \textbf{Mã HS:} \makebox[3.5cm]{\dotfill} \\
\hline
\end{tabularx}

\vspace{1.5mm}

\noindent\textbf{\large Câu 2 (2,5 điểm). Rút gọn và tìm $x$:}
\vspace{0.8mm}

\noindent
\begin{tabularx}{\linewidth}{|X|X|}
\hline
\oBaiTap[10]{a)}{Rút gọn và tính giá trị biểu thức:\par
$P = x^3 + 6x^2 + 12x + 8$\par
tại $x = 98$.} & 
\oBaiTap[10]{b)}{Tìm số thực $x$, biết:\par
$(x + 2)^3 - x^2(x + 6) = 44$.} \\
\hline
\end{tabularx}

\vspace{1.8mm}

\noindent\textbf{\large Câu 3 (2,5 điểm). Bài toán thực tế và nâng cao:}
\vspace{0.8mm}

\noindent
\begin{tabularx}{\linewidth}{|X|}
\hline
\vspace{0.5mm}
\textbf{a) (1,5 điểm)} Một bể nước hình lập phương có độ dài cạnh ban đầu là $x$ (mét). Người ta nâng cấp bằng cách tăng độ dài mỗi cạnh thêm $0{,}5\text{ m}$.\par
1) Hãy viết biểu thức biểu thị thể tích bể nước mới theo $x$. Viết biểu thức biểu thị lượng thể tích nước tăng thêm $\Delta V$ của bể mới so với bể cũ theo $x$.\par
2) Cho biết cạnh bể ban đầu là $x = 1\text{ m}$. Hãy tính chính xác lượng thể tích nước tăng thêm $\Delta V$ (theo $\text{m}^3$ và đổi ra lít nước, biết $1\text{ m}^3 = 1000$ lít).\par
\vspace{0.4mm}
\noindent\textit{Bài giải:}\dongke{7}
\vspace{1mm}
\textbf{b) (1,0 điểm)} Chứng minh rằng giá trị của biểu thức sau không phụ thuộc vào biến $x$:\par
$M = (x + 4)^3 - (x - 4)^3 - 24x^2$.\par
\vspace{0.4mm}
\noindent\textit{Bài giải:}\dongke{6}
\\ \hline
\end{tabularx}

\end{document}
